% ============================================================
%  preamble_beamer.tex
%  beamer 클래스 전용 공용 프리앰블 (XeLaTeX 기준)
%  beamer는 geometry/hyperref/fancyhdr을 자체적으로 관리하므로
%  여기서는 그것들을 별도로 불러오지 않는다.
% ============================================================

% ---------------------------------------------------------------
% 테마
% ---------------------------------------------------------------
\usetheme{Madrid}
\usecolortheme{default}
\setbeamertemplate{navigation symbols}{}
\setbeamertemplate{caption}[numbered]
\setbeamercovered{transparent}

% ---------------------------------------------------------------
% 한글 / 폰트 (XeLaTeX + kotex)
% ---------------------------------------------------------------
\usepackage{fontspec}
\usepackage{kotex}
% \setmainhangulfont{Noto Sans CJK KR}

% ---------------------------------------------------------------
% 수식 & 기호
% ---------------------------------------------------------------
\usepackage{amsmath, amssymb}
\usepackage{mathtools}
\usepackage{bm}
\usepackage{mathrsfs}
\usepackage{cancel}
\IfFileExists{physics.sty}{\usepackage{physics}}{}

\ifdefined\abs\else
  \DeclarePairedDelimiter{\abs}{\lvert}{\rvert}
\fi
\ifdefined\norm\else
  \DeclarePairedDelimiter{\norm}{\lVert}{\rVert}
\fi
\newcommand{\R}{\mathbb{R}}
\newcommand{\N}{\mathbb{N}}
\newcommand{\Z}{\mathbb{Z}}
\newcommand{\Q}{\mathbb{Q}}
\newcommand{\C}{\mathbb{C}}

% ---------------------------------------------------------------
% 표
% ---------------------------------------------------------------
\usepackage{booktabs}
\usepackage{multirow}
\usepackage{multicol}
\usepackage{tabularx}
\usepackage{makecell}

% ---------------------------------------------------------------
% 그림 & TikZ
% ---------------------------------------------------------------
\usepackage{xcolor}
\usepackage{tikz}
\usepackage{tikz-cd}
\usepackage{pgfplots}
\pgfplotsset{compat=1.18}
\usetikzlibrary{
  arrows.meta, calc, positioning, shapes.geometric, shapes.symbols,
  fit, backgrounds, patterns, matrix, intersections,
  decorations.pathmorphing, decorations.markings, angles, quotes
}
\tikzset{
  >=stealth,
  line cap=round,
  line join=round,
  box/.style={draw, rounded corners, minimum width=2.8cm, minimum height=1.2cm, align=center},
  point/.style={circle, fill, inner sep=1.2pt},
  arrow/.style={->, thick},
  highlight/.style={fill=blue!10, draw=blue!60!black}
}

% ---------------------------------------------------------------
% 코드 & 목록
% ---------------------------------------------------------------
\usepackage{listings}
\lstset{
  basicstyle=\ttfamily\footnotesize,
  breaklines=true,
  columns=fullflexible,
  frame=single,
}
